
% Authoritative LaTeX manuscript; Markdown is a generated compatibility projection.

\textbf{Problem.} Cross-cutting mechanisms depend on assumptions about the substrate they inspect. When those dependencies remain implicit, a substrate change can either break many controls at once or silently invalidate part of their coverage.

\textbf{Move.} Represent consequential dependencies explicitly so the impact of a proposed change can be queried before it lands.

\autoref{fig-move10} fans a change out to its dependents.

\begin{figure}[htbp]
\centering
\includegraphics[width=\linewidth]{book/latex/assets/c10-dependencies-queryable-c5b968db74.pdf}
\caption{\emph{Impact as a graph query.} A proposed change reaches a substrate node; typed dependency edges identify the controls and models that depend on it, and their union defines the affected set.}
\label{fig-move10}
\end{figure}

\textbf{Example --- Control blast radius.} Each control declares its substrate dependencies as structured metadata. A query identifies which controls depend on a proposed substrate change and how each relation should be interpreted.

\textbf{Example --- Traceability.} Symbol-anchored traceability applies the same move across engineering artifacts. Model elements, lints, code, proofs, and registries join through derived, re-checked edges. The resulting graph lets engineers query what a change affects and inspect why.

\textbf{Related mechanisms:} Control ↔ substrate dependency · Computed Control Blast Radius · Governance Graph · Symbol-anchored traceability · Derived Traceability.
